
\section*{初版にむけたあとがき}

このテキストは，広島大学総合科学部に非常勤講師としてお招きいただき，これまで筆者の学び，考えてきたことを「測定法」という観点からまとめ直したものになります。他の資料の流用・併用も多くありますが，新たに書き下ろしたセクションも少なくありません。心理測定，なかでも心理尺度に限った狭い話ではありますが，何かの参考になれば幸いです。

類書として，筆者の現所属での授業資料である「心理学データ解析基礎」，「心理学データ解析応用」のテキストや，数学ワークショップのためにまとめた「心理学者のための線形代数」などがあります。いずれも心理統計教育教材のサイト\footnote{\url{https://kosugitti.github.io/psychometrtics_syllabus/}}にてPDF他公開しておりますので，よろしければそちらもご参照ください。

これらのテキストを紙媒体でまとめたものも，実費で販売もしております。少々のお布施を厭わないという寛大な心をお持ちの方は，是非そちらもご利用ください。


なお原稿そのものはCreative Commons BY-SA(CC BY-SA)ライセンスVersion 4.0に基づいて提供しています。再頒布などは，ライセンスに従ってご利用いただければと思います。

\begin{flushright}
    \date{\today}\\
    小杉考司
\end{flushright}

